\Needspace{7\baselineskip}
\begin{lemma}[Packing blocks in a linear set system]\label{lem:block-capacity}
Let a set of $N$ points be partitioned into parts $\mathcal P$, and let
$\mathcal L$ be a family of subsets of size at least two. Distinct members
of $\mathcal L$ meet in at most one point, and each meets every part in
at most one point. For $q\ge1$, there are $q$ disjoint pairs, none inside a part and
at most one inside each member of $\mathcal L$, if and only if
\[
 N\ge2q,\qquad |P|\le N-q\ (P\in\mathcal P),\qquad
 |L|\le N-q+1\ (L\in\mathcal L).
\]
\end{lemma}
\begin{proof}
Necessity follows by counting points outside a part or a member of
$\mathcal L$. First suppose $N=2q$. For each $L$, add two adjacent
vertices and join both to every point of $L$. Join two original points
directly precisely when they lie in different parts and in no common
member of $\mathcal L$. A perfect matching of this graph gives the
required pairs: the two vertices of $L$ either match each other or
match two points of $L$. The intersection assumptions make these
assignments unambiguous.

Apply Tutte's theorem \cite{Tutte}. Delete a vertex set $S$, containing
$h$ original points, both added vertices of $r$ members of $\mathcal L$,
and one added vertex of $s$ others. Thus $|S|=h+2r+s$. Components with
no original point contribute at most $s$ odd components. If the number
$c$ of components containing points is at most $h+2r$, Tutte's inequality
follows. Assume $c>h+2r$.

If $c=1$, then $h=r=0$. For $s>0$, every remaining added vertex still
has a point neighbor, so there is at most one odd component, at most
$|S|$. For $s=0$, the graph has even order and its only point-containing
component is even. If $r=0$ and $c\ge2$, any two points in different
components must belong to the same part: a common member of
$\mathcal L$ would still have an added vertex joining them. All remaining points therefore
belong to one part, giving $c\le2q-h\le q\le h$, a contradiction.

It remains that $r\ge1$ and $c\ge3$. Choose one point from each component.
Every pair of representatives belongs to exactly one part or one of
the $r$ removed members of $\mathcal L$. If the representatives are not
all in one such set, the de Bruijn--Erd\H{o}s incidence theorem
\cite{deBruijnErdos} requires at least $c$ sets. There are at most
$r+\lfloor c/2\rfloor$, since the parts are disjoint. This contradicts
$c>h+2r$. Thus every choice of representatives is contained in one set.
Keeping two representatives and varying a third shows that all remaining
points lie in that same set. A part gives $c\le h$ as above. A member of
$\mathcal L$ gives $2q-h\le q+1$, hence
$c\le2q-h\le h+2\le h+2r$, again a contradiction. Tutte's theorem applies.

For general $N\ge2q$, add $N-2q$ dummy points as one new part, meeting
no member of $\mathcal L$. There are $2(N-q)$ points, and the balanced
case applies with $q$ replaced by $N-q$. Delete the pairs meeting
dummy points. The remaining $q$ pairs have the required properties.
\end{proof}

